\documentclass[11pt]{article}
\usepackage[T1]{fontenc}
\usepackage{lmodern,microtype}
\usepackage[margin=.94in]{geometry}
\usepackage{amsmath,amssymb,amsthm,booktabs,hyperref}
\newtheorem{theorem}{Theorem}
\newtheorem{lemma}[theorem]{Lemma}
\newcommand{\Z}{\mathbb Z}
\newcommand{\Q}{\mathbb Q}
\newcommand{\ee}{\varepsilon}
\DeclareMathOperator{\degop}{deg}
\title{Reciprocal numerators for shifted Hankel determinants of Catalan powers}
\author{Blake Stone's mathematics project\\ \small Developed with OpenAI Codex}
\date{November 26, 2026}
\setlength{\emergencystretch}{1.6em}
\begin{document}\maketitle
\begin{abstract}
For every positive Catalan power and every shift in the range of Cigler's generating-function conjectures, we prove a signed quasipolynomial formula for the Hankel determinants and determine the degree and reciprocity sign of the numerator. This proves Conjecture~15 of his 2023 preprint and gives Conjecture~16 with its single previously identified degree exception. For odd powers the resulting denominator bound improves the proposed exponent by the multiplicity of a zero of the moment polynomial. The proofs use a confluent determinant of Chebyshev polynomials, finite differences, and reflection in the determinant size.
\end{abstract}

\section{The generating functions}
For an integer $r\ge1$ let $C_{r,j}$ be the coefficient of $z^j$ in the $r$th power of the Catalan generating function. Equivalently,
\[
 C_{r,j}=\frac{r}{2j+r}\binom{2j+r}{j}\quad(j\ge0),
 \qquad C_{r,j}=0\quad(j<0).
\]
For an integer $m$ set $D(n)=D_{r,m}(n)=\det(C_{r,i+j+m})_{0\le i,j<n}$, with $D(0)=1$, and $F(z)=\sum_{n\ge0}D(n)z^n$.
Use the following parameters, where $k\ge1$, $m\ge1-k$, $b=k-1$, and $d=a+b$:
\[
\begin{array}{cccc}
\toprule
r & a & h & \ee\\\midrule
2k   & k+m   & k & (-1)^{\binom{k}{2}}\\
2k-1 & k+m-1 & 2k-1 & (-1)^{k-1}\\\bottomrule
\end{array}
\]
Write $e_0=1+\binom{a}{2}$, where $\binom{0}{2}=0$.
\begin{theorem}\label{thm}
For $0\le t<h$ there is a polynomial $p_t\in\Q[X]$ of degree at most $e_0-1$ such that
\[
 D(t+hj)=\ee^j p_t(j)\qquad(j\ge0).
\]
For each $e\ge e_0$, $P_e(z)=(1-\ee z^h)^e F(z)$ lies in $\Z[z]$. If $d\ge1$ its degree is $he-d$, its constant coefficient is $1$, and
\[
 z^{he-d}P_e(1/z)=-(-\ee)^e(-1)^{\binom{d}{2}}P_e(z).
\]
The sole case $d=0$ is $(r,m)=(1,0)$. Then $D(n)=1$ for all $n\ge0$ and $P_e(z)=(1-z)^{e-1}$.
\end{theorem}

For the even powers, set $e=e_0$. For the odd powers, Cigler's exponent in \cite[Conjectures 15--16, (28), (30)]{cigler} is $e_C=1+\binom{k+m}{2}=e_0+a$. In particular the odd-power numerator has the factor $(1-\ee z^h)^a$. Let \(\pi_s(t)=(s-2)\binom{t+1}{2}-t\). For $d\ge1$ Theorem~\ref{thm} gives the proposed degrees
\[
 \degop P_{e_0}=\pi_{k+2}(k+m-1)\quad(r=2k),\qquad
 \degop P_{e_C}=\pi_{2k+1}(k+m-1)+k\quad(r=2k-1).
\]
Substituting $a,d,\ee,e$ into the sign in the theorem gives the following residue classes of $m$ for which the source numerator is palindromic. The remaining classes give a skew-palindromic numerator. Entries list residues modulo four.
\[
\begin{array}{c|cccc}\toprule
 k\bmod4&0&1&2&3\\\midrule
 r=2k&1,3&0,1,2,3&0,3&1,2\\
 r=2k-1&0,1&0,2&0,1&1,3\\\bottomrule
\end{array}
\]
These are the sign classes of the two conjectures. At $(r,m)=(1,0)$ the source numerator $1$ is palindromic and its degree is zero. The literal degree prediction one at this parameter was already refuted in \cite{omega}. Wang and Xin~\cite{wang} studied the unshifted powers by continued fractions, including evaluations for $r\le31$. Our argument concerns the generating functions for all admissible shifts. It does not require conjectural evaluations of the determinants in any particular size class.

There is also overlap with the Catalan-like arrays of Chern and Shi~\cite{cs}. In their notation the weighted meander entry at height $k-1$ and length $u=j+k-1$ equals $C_{2k,j}$ at $(b,c)=(0,2)$, and equals $(-1)^j C_{2k-1,j}$ at $(b,c)=(1,-2)$. Indeed the even specialization is $\binom{2u}{u-k+1}-\binom{2u}{u-k-1}$ by reflection with two horizontal colors. For the odd specialization group ordinary nonnegative up-down steps in pairs; the horizontal multiplicities are one at height zero and two above. Negating both horizontal weights contributes $(-1)^{u-k+1}$. These specializations include the zero entries for $j<0$ whenever $u\ge0$. Their Theorems 1.1--1.9 give previous low-shift evaluations and individual size classes for variable shifts. Here the degree and reciprocity hold for the entire generating function at each admissible shift.

\section{A determinant in the size of the moment modification}
Let $\mu$ be the semicircle probability measure $\sqrt{4-y^2}\,dy/(2\pi)$ on $[-2,2]$ and put $L(f)=\int f(y^2)\,d\mu(y)$. Write $U_j$ for the Chebyshev polynomial of the second kind, with $U_j(\cos\theta)=\sin((j+1)\theta)/\sin\theta$. The substitution $y=2\cos\theta$ gives orthonormality of $U_j(y/2)$ for $\mu$. The monic orthonormal polynomials of $L$ are therefore
\[
 q_n(x)=U_{2n}(\sqrt{x}/2)
       =\sum_{i=0}^n(-1)^{n-i}\binom{n+i}{2i}x^i.
\]
Multiplication of $U_j(y/2)$ by $y$ sums its two neighbors, with $U_{-1}=0$. The coefficient of $U_{r-1}(y/2)$ in $y^N$ thus counts nonnegative nearest-neighbor paths of length $N$ from $0$ to $r-1$. For $N=2u+r-1\ge0$ reflection at $-1$ gives $\binom N u-\binom N{u-1}=C_{r,u}$ when $u\ge0$. For $u<0$ no path reaches $r-1$. Consequently
\[
 L(x^\ell M(x))=C_{r,\ell+m}\qquad(\ell\ge0),
 \quad M(y^2)=y^{2m+r-1}U_{r-1}(y/2).
\]
The exponent of $y$ is nonnegative in our shift range. The parity, leading coefficient and zeros of $U_{r-1}$ give
\[
 M(x)=x^a\prod_{j=1}^b(x-\alpha_j),\qquad
 \alpha_j=4\cos^2(j\pi/r).
\]
These $b$ interior roots are positive and distinct. The additional simple zero of $U_{r-1}$ at $y=0$ when $r$ is even is included in $x^a$.

Let $\Lambda$ be the ordered list of $d$ linear functionals consisting of coefficient extraction $[x^i]$ for $0\le i<a$, followed by evaluation at the $b$ roots $\alpha_j$. Define
\[
 E(n)=\big(\lambda(q_{n+j})\big)_{\lambda\in\Lambda,\,0\le j<d}.
\]
An empty determinant is one. The standard Christoffel formula for modified moments (compare \cite{kratt}) specializes as follows.
\begin{lemma}\label{chris}
For $n\ge0$, \(D(n)=(-1)^{nd}\det E(n)/\det E(0)\).
\end{lemma}
\begin{proof}
The evaluations $\Lambda$ are independent on polynomials of degree less than $d$, since their common kernel consists of multiples of $M$. Thus $\det E(0)\ne0$. On polynomials of degree less than $n+d$, use the ordered rows $L(x^\ell\,\cdot)$, $0\le\ell<n$, followed by $\Lambda$. In the monic basis $q_0,\ldots,q_{n+d-1}$ the top block is lower triangular with diagonal one in its first $n$ columns and vanishes in the remaining columns, by orthonormality. The full determinant is $\det E(n)$. Change to the ordered columns
\[
 M,xM,\ldots,x^{n-1}M,\quad 1,x,\ldots,x^{d-1}.
\]
Their degrees are a block permutation of $0,\ldots,n+d-1$ with sign $(-1)^{nd}$ and their leading coefficients are one. The lower left block vanishes, the upper left block has determinant $D(n)$, and the lower right block has determinant $\det E(0)$ by monicity. The stated identity follows.
\end{proof}

\section{Polynomial dependence and reflection}
For $N\in\Z$ use the evaluation formulas
\[
 [x^i]q_N=(-1)^{N-i}\binom{N+i}{2i},\qquad
 q_N(\alpha_j)=\frac{\sin((2N+1)j\pi/r)}{\sin(j\pi/r)}
\]
with the polynomial definition of the binomial coefficient. Both formulas also result from setting $q_{-N-1}=-q_N$ for $N\ge0$. Use Lemma~\ref{chris} to extend $D(n)$ to all integer $n$.

Expand $\det E(n)$ along the first $a$ rows. For selected columns with offsets $j_1,\ldots,j_a$ the corresponding endpoint minor, after extracting the signs, is
\[
 \det(B_i(n+j_\ell))_{0\le i<a,\,1\le\ell\le a},
 \qquad B_i(X)=\binom{X+i}{2i}.
\]
Its degree in $n$ is at most $\binom a2$. Indeed, use the finite Taylor expansion $B_i(n+j_\ell)=\sum_{v\ge0} B_i^{(v)}(n)j_\ell^v/v!$ columnwise. In the multilinear expansion, repeated derivative indices make two columns dependent. The sum of $a$ distinct nonnegative indices is at least $\binom a2$, and $\sum_{i=0}^{a-1}\degop B_i=a(a-1)$. This also covers empty minors.

On increasing $n$ by $h$, the endpoint signs change by $(-1)^{ha}$, while the prefactor of Lemma~\ref{chris} changes by $(-1)^{hd}$. For $r=2k$, each interior row $j$ changes by $(-1)^j$; for $r=2k-1$ every interior row stays fixed. The combined multipliers are respectively $(-1)^{k(k-1)+\binom{k}{2}}$ and $(-1)^{(2k-1)(k-1)}$. Hence on each residue the extended $D(n)$ is a polynomial of the asserted degree times the stated sign. Interpolation on nonnegative integer arguments gives rational coefficients. Finite differencing $e$ times on each signed arithmetic progression shows that $P_e$ is a polynomial, with integer coefficients from its defining series.

For $d\ge1$, the column list in $E(-n-d)$ is the reversed and negated list in $E(n)$. The factors from negating $d$ columns and from the two prefactors $(-1)^{nd}$ cancel because $d+d^2$ is even. Thus
\[
 D(-n-d)=(-1)^{\binom d2}D(n)\quad(n\in\Z).
\]
Moreover $D(-j)=0$ for $1\le j<d$, since those consecutive columns include both $q_{-1}=-q_0$ and $q_0$.

As rational functions expanded at zero we have
\[
 F(1/z)=-\sum_{j\ge1}D(-j)z^j.
\]
To see this directly, the bilateral sequence $D$ is annihilated by the convolution with $B(z)=(1-\ee z^h)^e$ by the preceding progression formula. Thus the positive half $\sum_{n\ge0}D(n)z^n$ and the negative half $\sum_{n<0}D(n)z^n$, multiplied separately by $B$, give opposite finite Laurent polynomials. Their rational functions sum to zero. Substitution of $1/z$ proves the displayed identity. The vanishing entries and reflection give $F(1/z)=-(-1)^{\binom d2}z^d F(z)$. Since $B(1/z)=(-\ee)^e z^{-he}B(z)$, this proves the stated reciprocal identity. Its constant term $P_e(0)=1$ forces degree exactly $he-d$. Finally, $d=0$ gives $M=1$ and $D(n)=1$ by Lemma~\ref{chris}.
\section{The explicit low shifts}
We give the six families displayed after (28) and (30) in \cite{cigler}. Write $s_j=(-1)^{\binom j2}$ for $j\ge0$. In this section $P$ is the numerator at the exponent in (28), and $Q$ the numerator at the exponent in (30). The even formulas are
\begin{align*}
 P_{2k,1-k}&=1 &&(k\ge1),\\
 P_{2k,2-k}&=1+s_{k-1}z^{k-1} &&(k\ge2),\\
 P_{2k,3-k}&=(1+s_{k-2}z^{k-2})(1-z^{2k})\\
 &\quad+s_{k-1}k^2z^{k-1}(1+s_k z^k) &&(k\ge3).
\end{align*}
Put $U=1+(-1)^kz^{2k-1}$. For the odd powers,
\begin{align*}
 Q_{2k-1,1-k}&=1+s_k z^k &&(k\ge2),\\
 Q_{2k-1,2-k}&=(1+s_{k-1}z^{k-1})U &&(k\ge2),\\
 Q_{2k-1,3-k}&=U^2\big[(1+s_{k-2}z^{k-2})U\\
 &\qquad\quad +(2k-1)z^{k-1}(z^{k-1}+s_{k-1})\big] &&(k\ge3).
\end{align*}
The last bracket is precisely the expression printed on p.~11 of \cite{cigler}. It uses the exponent $e_0=2$; equation (30) calls for $e_C=4$, hence the external factor $U^2$. Also $C_{2,j}=C_{1,j+1}$ for $j\ge0$, so for $m\ge0$ the comparison after (30) reads $Q_{1,m+1}=(1-z)^{m+1}P_{2,m}$. Reversing $1-z$ changes that identity's sign when $m$ is even.

For the first two shifts in each parity class, Theorem~\ref{thm} gives the degree of the numerator at $e_0$. If $m=-M\le0$, then $D(n)=0$ for $1\le n\le M$ by the first row, and $D(M+1)=s_{M+1}$ by the opposite diagonal of ones. The stated low-shift identities in these four cases follow from the degree and the initial coefficients, inserting $U^{e_C-e_0}$ for the odd powers.

For $m=3-k$ put $\ell=k-2$. Write $A(z)=C(z)^r$ and $B(z)=1/A(z)=\sum b_j z^j$, where $C$ is the Catalan series. For $N=\ell+q$, $q\ge0$, complementary minors give
\begin{equation}\label{comp}
 D(N)=s_\ell(-1)^q\det(b_{\ell+1+i+j})_{0\le i,j<q}.
\end{equation}
Indeed, in the $(N+q)$-dimensional unit lower triangular Toeplitz matrix of $A$, take rows $q,\ldots,q+N-1$ and columns $0,\ldots,N-1$. Reversing its columns gives the Hankel determinant up to $s_N$. The inverse is the Toeplitz matrix of $B$. Jacobi's complementary minor identity gives the complementary inverse minor with rows $N,\ldots,N+q-1$, columns $0,\ldots,q-1$, and sign $(-1)^{Nq}$. Reverse its columns to obtain the indicated $b$-determinant. The total sign $s_Ns_q(-1)^{Nq}=s_{\ell+2q}=s_\ell(-1)^q$.

The roots $\lambda_\pm=(1\pm\sqrt{1-4z})/2$ satisfy $\lambda_+=C(z)^{-1}$ and $\lambda_-=zC(z)$. Hence $B(z)+z^r A(z)=\lambda_+^r+\lambda_-^r$. By the recursion $T_r=T_{r-1}-zT_{r-2}$ with $T_0=2,T_1=1$, the latter polynomial has coefficients
\[
 [z^j]T_r=(-1)^j\frac{r}{r-j}\binom{r-j}{j},\quad
 0\le j\le\lfloor r/2\rfloor .
\]
For $k-1\le j<r$, the nonzero $b_j$ are therefore
\[
 \begin{array}{c|cc}
   & b_{k-1}& b_k\\\hline
 r=2k   &(-1)^{k-1}k^2&2(-1)^k\\
 r=2k-1 &(-1)^{k-1}(2k-1)&0 .
 \end{array}
\]
For even $r$, the degree is $3k-2$, so reciprocity only needs coefficients through $N=\lfloor(3k-2)/2\rfloor$. For odd $r$ use $e_0$ first; its degree is $3k-3$ so $N\le\lfloor(3k-3)/2\rfloor$ suffices. Throughout these ranges the largest $b$-index in \eqref{comp} is $\ell+2q-1=2N-\ell-1<r$. The right-hand determinant vanishes for $q\ge3$ in the even case and $q\ge2$ in the odd case, by its last row. At $q=0$ it is one. At $q=1$ it is $b_{k-1}$, and at $q=2$ in the even case it is $-b_k^2=-4$.

In both cases the lower half of the numerator thus consists of coefficient one at zero, $s_{k-2}$ at $k-2$, and $s_{k-1}k^2$ (even) or $s_{k-1}(2k-1)$ (odd) at $k-1$. In the even case the intermediate $D(k)=-4s_{k-2}$ cancels the contribution $-4s_kD(0)$ from the multiplier since $s_k=-s_{k-2}$. In the odd case the multiplier starts beyond the lower half. Theorem~\ref{thm} reflects these three coefficients with sign $s_k$ for the even case and $s_{k-1}$ for the smaller odd numerator. This proves the two remaining expressions.
\begin{thebibliography}{9}
\bibitem{cigler} J. Cigler, \emph{Some experimental observations about Hankel determinants of convolution powers of Catalan numbers}, v2, 2023, \S3, Conjectures 15 and 16. \href{https://arxiv.org/abs/2308.07642v2}{arXiv:2308.07642v2}.
\bibitem{cs} S. Chern and W. Shi, \emph{Hankel determinants of Catalan-like sequences} (2026). \href{https://arxiv.org/abs/2608.27208}{arXiv:2608.27208}.
\bibitem{kratt} C. Krattenthaler, \emph{Hankel determinants of linear combinations of moments of orthogonal polynomials, II}, Ramanujan J. \href{https://doi.org/10.1007/s11139-021-00514-8}{doi:10.1007/s11139-021-00514-8}.
\bibitem{omega} TrureTuring, \emph{The literal numerator degree in Cigler's Conjecture 16}, research record for $(k,m)=(1,0)$. \href{https://www.trureturing.org/research/cigler-sixteen-literal-numerator-degree/}{Public record}.
\bibitem{wang} Y. Wang and G. Xin, \emph{Hankel determinants for convolution powers of Catalan numbers}, Discrete Math. 342 (2019), 2694--2716. \href{https://doi.org/10.1016/j.disc.2019.05.025}{doi:10.1016/j.disc.2019.05.025}.
\end{thebibliography}
\end{document}
